\documentclass[11pt,a4paper]{article}
\usepackage[utf8]{inputenc}
\usepackage{amsmath,amssymb}
\usepackage{geometry}
\geometry{margin=2.5cm}
\usepackage{listings}
\usepackage{xcolor}
\usepackage{hyperref}

\lstset{
  language=C++,
  basicstyle=\ttfamily\small,
  keywordstyle=\color{blue}\bfseries,
  commentstyle=\color{green!60!black},
  stringstyle=\color{red!70!black},
  numbers=left,
  numberstyle=\tiny\color{gray},
  numbersep=8pt,
  frame=single,
  breaklines=true,
  tabsize=4,
  showstringspaces=false
}

\title{IOI 1998 -- Camelot}
\author{Solution}
\date{}

\begin{document}
\maketitle

\section{Problem Statement}

On an $R \times C$ chessboard ($R \le 30$, $C \le 26$), there is one king and
up to 63 knights. All pieces must gather at a single square. A knight may optionally
pick up the king en route: it visits the king's square, carries the king, and
continues to the gathering point. The king walks using Chebyshev distance (one step
in any of 8 directions per move).

Find the minimum total number of moves for all pieces to gather.

\section{Solution Approach}

\subsection{BFS for Knight Distances}

Precompute the shortest knight distance between every pair of squares using BFS from
each square. Let $\mathrm{kdist}[s][t]$ denote this distance. The board has at most
$S = 30 \times 26 = 780$ squares.

\subsection{King Distance}

The king's distance is the Chebyshev distance:
\[
  \mathrm{kingDist}((r_k, c_k), (r, c)) = \max(|r_k - r|, |c_k - c|).
\]

\subsection{Enumeration}

For each potential gathering square $g$, compute:

\paragraph{Option 1 -- King walks alone.}
\[
  \mathrm{cost}_1 = \sum_{i} \mathrm{kdist}[\mathrm{knight}_i][g] + \mathrm{kingDist}(\mathrm{king}, g).
\]

\paragraph{Option 2 -- A knight picks up the king.}
Knight $i$ detours through pickup square $p$, where the king walks to $p$:
\[
  \mathrm{cost}_2 = \sum_{j \ne i} \mathrm{kdist}[\mathrm{knight}_j][g]
    + \mathrm{kdist}[\mathrm{knight}_i][p] + \mathrm{kdist}[p][g]
    + \mathrm{kingDist}(\mathrm{king}, p).
\]

Take the minimum over all $g$, all knights $i$, and all pickup points $p$.

\subsection{Pickup Point Search}

To guarantee correctness, we search \textbf{all} squares as potential pickup points.
Although this increases the time to $O(S^2 \cdot n)$, it remains feasible for the
given constraints ($S \le 780$, $n \le 63$).

\section{C++ Solution}

\begin{lstlisting}
#include <cstdio>
#include <cstring>
#include <queue>
#include <algorithm>
#include <climits>
using namespace std;

const int MAXR = 30, MAXC = 26;
int R, C;
int kdist[MAXR * MAXC][MAXR * MAXC];
int dx[] = {-2, -2, -1, -1, 1, 1, 2, 2};
int dy[] = {-1, 1, -2, 2, -2, 2, -1, 1};

int idx(int r, int c) { return r * C + c; }

void bfs(int src) {
    queue<int> q;
    q.push(src);
    kdist[src][src] = 0;
    while (!q.empty()) {
        int u = q.front(); q.pop();
        int r = u / C, c = u % C;
        for (int d = 0; d < 8; d++) {
            int nr = r + dx[d], nc = c + dy[d];
            if (nr < 0 || nr >= R || nc < 0 || nc >= C) continue;
            int v = idx(nr, nc);
            if (kdist[src][v] == -1) {
                kdist[src][v] = kdist[src][u] + 1;
                q.push(v);
            }
        }
    }
}

int main() {
    scanf("%d %d", &R, &C);
    int total = R * C;

    // Read king position
    char col;
    int row;
    scanf(" %c %d", &col, &row);
    int kingR = row - 1, kingC = col - 'A';

    // Read knight positions
    int knights[65], numKnights = 0;
    while (scanf(" %c %d", &col, &row) == 2) {
        int kr = row - 1, kc = col - 'A';
        knights[numKnights++] = idx(kr, kc);
    }

    // BFS from every square
    memset(kdist, -1, sizeof(kdist));
    for (int i = 0; i < total; i++)
        bfs(i);

    // Replace unreachable entries with a large sentinel
    const int INF = 1000000;
    for (int i = 0; i < total; i++)
        for (int j = 0; j < total; j++)
            if (kdist[i][j] == -1)
                kdist[i][j] = INF;

    // Handle edge case: no knights
    if (numKnights == 0) {
        printf("0\n");
        return 0;
    }

    int ans = INT_MAX;

    // Try each gathering point
    for (int g = 0; g < total; g++) {
        int gr = g / C, gc = g % C;

        // Sum of all knight distances to g
        int base = 0;
        bool reachable = true;
        for (int i = 0; i < numKnights; i++) {
            if (kdist[knights[i]][g] >= INF) { reachable = false; break; }
            base += kdist[knights[i]][g];
        }
        if (!reachable) continue;

        // Option 1: king walks alone to g
        int kingAlone = base + max(abs(kingR - gr), abs(kingC - gc));
        ans = min(ans, kingAlone);

        // Option 2: one knight picks up king at some pickup point p
        for (int i = 0; i < numKnights; i++) {
            int saved = kdist[knights[i]][g];

            for (int p = 0; p < total; p++) {
                int pr = p / C, pc = p % C;
                int kingWalk = max(abs(kingR - pr), abs(kingC - pc));
                int knightDetour = kdist[knights[i]][p] + kdist[p][g];
                if (knightDetour >= INF) continue;
                int cost = base - saved + knightDetour + kingWalk;
                ans = min(ans, cost);
            }
        }
    }

    printf("%d\n", ans);
    return 0;
}
\end{lstlisting}

\section{Complexity Analysis}

\begin{itemize}
  \item \textbf{Time complexity:}
    \begin{itemize}
      \item BFS from every square: $O(S^2)$ where $S = R \times C \le 780$.
      \item Main enumeration: for each gathering point ($O(S)$), for each knight ($O(n)$),
            for each pickup point ($O(S)$): $O(S^2 \cdot n)$.
      \item Total: $O(S^2 \cdot n) \approx 780^2 \times 63 \approx 38{,}000{,}000$.
    \end{itemize}

  \item \textbf{Space complexity:} $O(S^2)$ for the knight distance table ($\approx 2.4$ MB).
\end{itemize}

\end{document}
